\chapter{Introducción, contexto, alcance, objetivos y metodología}\label{cap:intro}

\section{Introducción}\label{sec:introduccion}

La cosmología moderna ---la expansión del universo, la ley de Hubble-Lemaître, la determinación de la edad y el tamaño del universo y, en su frontera, la energía oscura y la expansión acelerada--- es uno de los contenidos científicos que mayor fascinación despierta en el alumnado y, al mismo tiempo, uno de los más expuestos a ideas erróneas persistentes y a un tratamiento escolar superficial. La literatura educativa muestra que estos temas motivan de forma consistente, pero que esa motivación no se traduce automáticamente en comprensión conceptual, especialmente cuando se recurre a analogías sin explicitar sus límites o cuando la energía oscura se presenta como una ``sustancia conocida'' en lugar de como un concepto introducido para dar cuenta de una observación.

Este Trabajo de Fin de Máster (TFM), enmarcado en el Máster Universitario en Astronomía y Astrofísica de la Universidad Internacional de Valencia (VIU), propone el \textbf{diseño de una secuencia didáctica basada en el modelo 5E} (\textit{Engage, Explore, Explain, Elaborate, Evaluate}) con un objetivo concreto: que el alumnado de 2º de Bachillerato llegue a \textbf{determinar la edad y el tamaño del universo}. La secuencia, \textbf{basada en un problema}, parte de una pregunta motriz ---¿podemos determinar la edad y el tamaño del universo y cómo lo sabe la ciencia?--- y comienza por \textbf{clarificar y redefinir los conceptos implicados} ---el espacio, el tiempo, el universo y qué se entiende por su «edad» y su «tamaño»--- y avanza mostrando \textbf{cómo, a partir de observaciones} ---el corrimiento al rojo de las galaxias y la ley de Hubble-Lemaître---, esas magnitudes pueden medirse e inferirse. La propuesta se fundamenta en una revisión sistematizada de la literatura y se ancla en el currículo de Física real del contexto de implementación.

\section{Contexto}\label{sec:contexto}

\textbf{Contexto profesional.} El autor ejerce como profesor de Física (y de \textit{Computer Science}) en educación secundaria y bachillerato. Imparte \textbf{Física de Pearson Edexcel International Advanced Level (IAL)} en 1º y 2º de Bachillerato y prepara a alumnado para la \textbf{Prueba de Competencias Específicas (PCE) de Física de la UNED} (2º de Bachillerato, currículo LOMLOE --- RD 243/2022). Este doble marco curricular es determinante para el trabajo: como se detalla en el Capítulo 2, los contenidos objeto del TFM constituyen \textbf{currículo evaluable} en la \textbf{Unit~5.6 ``Astrophysics and Cosmology''} de Edexcel (Año 2; en el libro de texto, subtema \textit{11B ``Space''} del \textit{Topic~11}) y aparecen, de forma más repartida, en los Bloques II y VI de la PCE. El autor dispone, por tanto, de un escenario natural para una futura implementación en aula.

\textbf{Contexto educativo del problema.} Pese a su presencia curricular, la enseñanza de la cosmología moderna afronta barreras bien documentadas: déficit de formación específica del profesorado, saturación curricular, escaso tiempo lectivo y ausencia de secuencias didácticas fundamentadas y contrastadas. El resultado es que un contenido de alto potencial motivador se imparte con frecuencia de forma expositiva y superficial.

\section{Justificación}\label{sec:justificacion}

El trabajo se justifica por la convergencia de cuatro factores: \textbf{(i)} el fuerte poder motivador de la cosmología, confirmado por la literatura; \textbf{(ii)} la existencia de un \textbf{núcleo estable de ideas erróneas} que requiere un tratamiento didáctico dirigido; \textbf{(iii)} el encaje directo del tema en el currículo que imparte el autor, lo que neutraliza la objeción de ``saturación curricular'' (la secuencia enseña contenido ya prescrito); y \textbf{(iv)} un hueco identificado en la literatura: la \textbf{determinación de la edad y el tamaño del universo} apenas se trabaja de forma explícita como objetivo de aprendizaje, lo que abre un margen de originalidad para la propuesta.

\section{Alcance y delimitación}\label{sec:alcance}

\begin{itemize}
\item \textbf{Naturaleza del trabajo.} Se trata de un TFM de \textbf{diseño}: su producto es una secuencia didáctica 5E fundamentada y sus instrumentos de evaluación. \textbf{No incluye implementación ni validación empírica} dentro del plazo del trabajo (redacción prevista de julio a septiembre de 2026). El autor asume el compromiso de \textbf{implementar la propuesta durante el curso 2026-2027} como fase posterior, si la línea de trabajo continúa.
\item \textbf{Nivel educativo.} 2º de Bachillerato (equivalente al Año 2 de Edexcel IAL; edad \textasciitilde17--18 años).
\item \textbf{Delimitación temática.} Núcleo: la determinación de la edad y el tamaño del universo, a través de la ley de Hubble-Lemaître y la expansión del espacio. Matiz (fase Elaborate): expansión acelerada y energía oscura, como límite de la estimación. Contenidos de apoyo: corrimiento al rojo, radiación cósmica de fondo, materia oscura. \textbf{Fuera del alcance}: el formalismo de la relatividad general y las funciones de distancia cosmológicas, reservados a alumnado avanzado.
\item \textbf{Contexto curricular de referencia.} Edexcel IAL Physics (Unit~5.6; en el libro, subtema 11B ``Space'' del \textit{Topic~11}) y PCE UNED de Física (LOMLOE).
\end{itemize}

\section{Objetivos}\label{sec:objetivos}

\subsection{Objetivo general}

Diseñar una secuencia didáctica basada en el modelo 5E para la determinación de la edad y el tamaño del universo ---empleando la expansión y la ley de Hubble-Lemaître como vía de medida, y con la expansión acelerada y la energía oscura como matiz--- en 2º de Bachillerato, fundamentada en una revisión sistematizada de la literatura y coherente con el currículo de Física del contexto de implementación. La propuesta incorpora, además, sus instrumentos de evaluación y una metodología para la lectura e interpretación de las ecuaciones que sustentan la medida.

\subsection{Objetivos específicos}

\begin{itemize}
\item \textbf{OE1.} Realizar una revisión sistematizada (\textit{scoping review}) de la literatura sobre enseñanza y divulgación de cosmología moderna en secundaria, identificando los contenidos más tratados, las estrategias empleadas y las dificultades conceptuales documentadas \textit{(Fase 1 $\rightarrow$ Cap. 2)}.
\item \textbf{OE2.} Categorizar las estrategias y extraer los elementos transferibles, asignándolos a las fases del modelo 5E \textit{(Fase 1 $\rightarrow$ Cap. 2)}.
\item \textbf{OE3.} Cotejar los contenidos de la revisión con el currículo de Física que imparte el autor (Edexcel IAL y PCE UNED) para fijar el nivel y el encaje curricular de la propuesta \textit{(Fase 1 $\rightarrow$ Cap. 2)}.
\item \textbf{OE4.} Diseñar la secuencia 5E seleccionando, a partir de la síntesis, qué contenidos incluir, qué estrategias son más adecuadas y qué ideas erróneas deben abordarse de forma explícita \textit{(Fase 2 $\rightarrow$ Cap. 3)}.
\item \textbf{OE5.} Diseñar los instrumentos de evaluación: un \textbf{cuestionario conceptual pre-post} (adaptado de inventarios validados) y \textbf{dos rúbricas} ---del trabajo de indagación y del producto final---, previendo su aplicación en la implementación del curso 2026-2027 \textit{(Fase 2 $\rightarrow$ Cap. 4)}.
\item \textbf{OE6.} Proponer una metodología didáctica para la lectura e interpretación de las ecuaciones que sustentan la medida (no su deducción formal), como aportación transversal a la secuencia \textit{(Fase 2 $\rightarrow$ Cap. 4)}.
\end{itemize}

\section{Metodología general}\label{sec:metodologia}

El TFM adopta un enfoque de \textbf{investigación educativa basada en el diseño} (\textit{design-based research}), en su \textbf{fase de diseño}: la propuesta se fundamenta empíricamente en la literatura y en el análisis curricular, y su validación en aula queda diferida a una fase posterior. El trabajo se organiza en dos fases metodológicas.

\subsection{Fase 1 --- Revisión bibliográfica y análisis (Capítulo 2)}

Constituye el grueso del trabajo realizado y la base documental que sustenta la propuesta. Comprende: una búsqueda semántica amplia (Elicit Pro), un cribado por criterios (\textit{screening}), una extracción de datos en diez dimensiones, la depuración y categorización del corpus, y la síntesis en torno a tres preguntas de revisión (contenidos, estrategias y dificultades), rematada con el cotejo curricular. Todo el procedimiento se describe y justifica en el Capítulo 2, de modo que resulte trazable cómo se construyó la base documental.

\subsection{Fase 2 --- Diseño y desarrollo de la secuencia didáctica 5E (Capítulos 3 y 4)}

A partir de la síntesis de la Fase 1 se diseñará la propuesta, organizada en las fases \textit{Engage, Explore, Explain, Elaborate} y \textit{Evaluate}. Siguiendo la orientación del director, los resultados de la revisión se emplearán para decidir de forma fundamentada:

\begin{itemize}
\item \textbf{qué contenidos incluir} en cada fase (a partir de PR1 y del anclaje curricular);
\item \textbf{qué estrategias son más adecuadas} (a partir de la categorización de PR2 y de la discusión);
\item \textbf{qué dificultades conceptuales deben abordarse} de forma explícita (a partir de PR3); y
\item \textbf{cómo evaluar la comprensión} del alumnado (a partir de los instrumentos de evaluación y del mapeo 5E).
\end{itemize}

El resultado se articula en dos capítulos: su \textbf{diseño} (Capítulo 3) ---objetivos, fases, encaje curricular y estructura--- y su \textbf{desarrollo} (Capítulo 4) ---actividades paso a paso, materiales e instrumentos de evaluación, para los que se propone una web interactiva como soporte, y una metodología para leer e interpretar las ecuaciones que sustentan la medida---.

\subsection{Cronograma}

\begin{table}[ht]
\centering
\begin{tabularx}{\textwidth}{|>{\raggedright\arraybackslash}p{2.6cm}|>{\raggedright\arraybackslash}X|>{\raggedright\arraybackslash}p{3.6cm}|} \hline
\rowcolor{naranja}
{\color{white}\textbf{Fase}} & {\color{white}\textbf{Contenido}} & {\color{white}\textbf{Periodo}} \\ \hline
Fase 1 & Revisión bibliográfica, análisis y anclaje curricular & Realizada / consolidación jul. 2026 \\ \hline
Fase 2 & Diseño de la secuencia 5E e instrumentos & Jul.--sep. 2026 \\ \hline
Redacción del TFM & --- & Jul.--sep. 2026 \\ \hline
Depósito del TFM & --- & Septiembre u octubre de 2026 (por decidir) \\ \hline
(Post-TFM) Implementación en aula y validación & Puesta en práctica y análisis pre-post & Curso 2026-2027 \\ \hline
\end{tabularx}
\end{table}

\section{Estructura del documento}\label{sec:estructura}

\begin{itemize}
\item \textbf{Capítulo 1.} Introducción, contexto, alcance, objetivos y metodología \textit{(este capítulo)}.
\item \textbf{Capítulo 2.} Revisión bibliográfica: construcción de la base documental \textit{(desarrollo de la Fase 1)}.
\item \textbf{Capítulo 3.} Diseño de la secuencia didáctica 5E \textit{(Fase 2: diseño)}.
\item \textbf{Capítulo 4.} Desarrollo de la secuencia didáctica 5E: actividades, materiales e instrumentos, con una web interactiva como soporte propuesto \textit{(Fase 2: desarrollo)}.
\item \textbf{Capítulo 5.} Conclusiones y líneas futuras (incluida la implementación 2026-2027).
\item \textbf{Anexo A.} Fuentes de la revisión exploratoria (Elicit Pro): el listado de fuentes identificadas, al que remiten los índices [n] de la tabla de análisis del Capítulo 2.
\item \textbf{Anexo B.} Las ecuaciones de la cosmología: la lectura, según la metodología presentada en el §4.8, de las nueve ecuaciones que vertebran la medida ---las cinco que el alumnado emplea en la secuencia y las cuatro de fondo, procedentes de la relatividad general, que se leen sin derivarlas.
\end{itemize}